\begin{frame}{Définition}
\begin{block}{}
    Ensemble de règles qui permettent à deux équipements informatiques (ordinateurs, smartphones, serveurs) d'échanger des données de manière compréhensible et ordonnée.
\end{block}

\begin{block}{Les protocoles de l'Internet et du Web}
    \begin{itemize}
        \item TCP (\textit{Transmission Control Protocol})
        \item IP (\textit{Internet Protocol})
        \item HTTP (\textit{Hypertext Transfer Protocol})
    \end{itemize}
\end{block}
    
\end{frame}

\begin{frame}{Le protocole HTTP}
    \begin{block}{}
        \centering
        {\small \textit{Le web, ou plus précisément le World Wide Web, est un système documentaire construit sur Internet dans lequel les documents, nommés hypertextes ou pages web, sont reliés les uns aux autres par des hyperliens.Ces documents sont affichés dans des navigateurs qui permettent, grâce aux hyperliens, de naviguer d’une page à une autre. Les données qui sont échangées sur le web le sont avec le protocole HTTP, HyperText Transfert Protocol.}}\\
        \textbf{Jean-Philippe Magué, "Les protocoles d’Internet et du web", Pratiques de l'édition numérique, PUM, 2024}
    \end{block}

    \begin{block}{}
        \textrightarrow{} Fonctionne sur le principe de requête et réponse.
    \end{block}
\end{frame}

\begin{frame}{Apprendre à lire une url}
\begin{block}{}
    \textbf{URL} = protocole, adresse d'un serveur, emplacement et le nom d'une ressource \\
    \textbf{Ex}: \url{https://fr.wikipedia.org/wiki/Wikipédia:Pastiches/Une_femme}
\end{block}
\begin{block}{}
    \begin{itemize}
        \item https:// = le protocole
        \item fr.wikipedia.org = le nom de domaine
        \item /wiki/Wikipédia:Pastiches/Une\_femme = le chemin
    \end{itemize}
\end{block}

\begin{block}{Structure interne de la page Web}
    \begin{itemize}
        \item wiki = le nom de l'outil d'édition
        \item wikipédia: = espace de noms. Signifie que nous sommes dans les pages communautaires ou d'administration du site.
        \item Pastiches = sous-catégorie ou dossier regroupant des textes humoristiques
        \item Une\_femme = titre de la page.
    \end{itemize}
\end{block}
    
\end{frame}

\begin{frame}{Les protocoles éditoriaux de l'Internet : une émanation de la contre-culture américaine}

\begin{block}{Fondements d'une contre-culture basée sur un idéal d'accessibilité des ressources et de liberté d'échange des contenus.}
\vspace{1em}
\begin{itemize}
    \item URL/URI (Tim Berners-Lee): principe d'adressage universel de toute information. Intégration des notions d'accès, partage ou appropriation des outils de publication.
    \item HTTP permet d'accéder à une ressource sans demander la permission au propriétaire du serveur.
    \item Tout média, toute information, toute donnée est susceptible sur le web d’être identifiée, extraite, associée, partagée, détournée.
\end{itemize}
    \textrightarrow{}  Culture de la libre circulation
\end{block}

Ces "règles" sont constamment négociées (par les grands groupes, les GAFAM, les chercheurs, la société civile, etc.)

\end{frame}

\begin{frame}{Exemple de travaux de recherche sur cette contre-culture}
    \centering
    {\small
    \textit{Les hippies ont projeté leur rêve d’exil et de refondation dans les échanges numériques et, pour cela, ils avaient besoin de couper les ponts avec un « réel » doublement décevant, en raison de la persistance de l’aliénation patriarcale et capitaliste, mais aussi de l’échec de la tentative de s’en émanciper en établissant dans ses marges des communautés contre-culturelles. Internet était un « ailleurs », le nouvel asile d’un projet d’émancipation avorté.}}\\
    \textbf{Fred Turner, \textit{Aux sources de l’utopie numérique}, C\&F édition, 2021 (2006)}
\end{frame}

\begin{frame}{Exemple de travaux de recherche sur cette contre-culture}
    \centering
    {\small \textit{Néanmoins, la rhétorique de « l’informationalisme en pair à pair », rejoignant en cela la rhétorique de la conscience dont elle était issue, occulte fortement les infrastructures matérielles et techniques dont dépendent à la fois l’Internet et les vies des membres de la génération numérique. Derrière le fantasme d’un flux d’information sans contrainte, s’étant la réalité de millions de claviers en plastique, de puces en silicium, de moniteurs aux écrans de verre, et de kilomètres de câbles à l’infini.}}\\
    \textbf{Fred Turner, \textit{Aux sources de l’utopie numérique}, C\&F édition, 2021 (2006)}
\end{frame}